% ---------------------------------------------------------------------
%  Chapitre 2 — Les données et leur structure de dépendance
% ---------------------------------------------------------------------
\chapter{Mesurer la dépendance qui menace le jugement}
\label{chap:donnees}

\section{La base des trades}
\label{sec:donnees-trades}


Tout au long de ce rapport, on appliquera les méthodes à notre base
de données. Celle-ci a été obtenue en appliquant différentes stratégies
sur des données antérieures. Cette méthode se nomme le \emph{backtesting}, et
consiste à vérifier si notre stratégie a repéré un signal de prédiction
du marché.

Ici nous appliquons 10 stratégies techniques classiques ainsi qu'une stratégie
oracle témoin, soit onze stratégies au total.
La stratégie oracle utilise les données futures pour effectuer des
trades uniquement gagnants, afin de vérifier si notre protocole laisse passer
une stratégie dont le signal fonctionne réellement.
Les 10 autres stratégies, elles, effectuent un trade à une date fixée en
utilisant uniquement les données précédant cette date. 

\begin{table}[H]
\centering
\caption{Les onze stratégies du banc d'essai et l'indicateur sur lequel chacune repose.}
\label{tab:strategies}
\renewcommand{\arraystretch}{1.25}
\small
\begin{tabular}{@{}ll@{}}
\toprule
Stratégie & Fondée sur \\
\midrule
\texttt{rsi\_classic}  & indicateur RSI (seuils 30/70) \\
\texttt{rsi\_strict}   & indicateur RSI (seuils 20/80) \\
\texttt{rsi\_trend}    & indicateur RSI avec filtre de tendance (MM200) \\
\texttt{ma\_crossover} & croisement de deux moyennes mobiles \\
\texttt{db\_bottom}    & figure en double creux \\
\texttt{dt\_top}       & figure en double sommet \\
\texttt{hs\_inverse}   & figure épaule-tête-épaule inversée \\
\texttt{hs\_classic}   & figure épaule-tête-épaule \\
\texttt{sr\_breakout}  & cassure d'une résistance \\
\texttt{sr\_breakdown} & cassure d'un support \\
\midrule
\texttt{oracle}        & témoin : connaît le cours futur (triche) \\
\bottomrule
\end{tabular}
\end{table}

J'ai alors choisi d'appliquer ces stratégies de trading aux grandes
capitalisations américaines présentes dans le S\&P~500. Cet indice suit les
500 plus grandes sociétés cotées aux États-Unis, ce qui représente ici
501 titres, quelques sociétés étant cotées avec deux catégories d'actions.
C'est un ensemble sur lequel on peut investir avec peu de risque : même si le
cours de chacune de ces actions peut chuter individuellement, l'indice
lui-même reste relativement stable et possède une tendance haussière.

Pour qu'une stratégie soit validée, il faut qu'elle ait de meilleurs
résultats qu'une stratégie au hasard et que cela ne soit pas dû à
notre échantillonnage.
Ici la tendance est haussière, ce qui nous amène à la première
difficulté mentionnée dans l'introduction : comment évaluer une 
stratégie lorsque la tendance est haussière ?

Pour répondre à cette problématique, on va simuler le hasard.
À chaque trade réalisé, on associe 200~trades fictifs de même exposition
(même titre, même durée, même sens), mais dont la date
d'entrée est tirée au sort ; la moyenne de leurs rendements définit le
gain du hasard pour ce trade.

On définit alors, pour chaque trade, son \edgevar{} comme le rendement
effectivement obtenu par la stratégie (éventuellement négatif) diminué de
ce gain du hasard :
\[
  \edgevar = \text{gain de la stratégie} - \text{gain du hasard}.
\]
Un \edgevar{} positif signifie que la stratégie a fait mieux qu'une entrée
au hasard de même exposition : son \emph{timing} apporte une information
que la seule dérive du marché n'explique pas.

On obtient ainsi une base de données contenant, pour chaque stratégie,
entre 978 et 8\,513 trades effectués entre janvier~2010 et juillet~2026,
pour un total de 45\,008 trades.
Dans nos données, chaque ligne correspond à un trade, c'est-à-dire une
position ouverte puis fermée par une stratégie sur un titre : la stratégie
achète l'action à une date donnée, la conserve un certain temps, puis la
revend. C'est cette ligne qui constituera notre unité statistique dans tout
le reste du rapport.
On y associe les indications suivantes : le titre
(\emph{ticker}), la stratégie employée, le secteur, les dates d'entrée et de
sortie, la durée de détention, le régime de marché et le contexte à l'entrée
(volatilité, RSI, distance à la moyenne mobile), le rendement et l'avantage
(\edgevar{}) sur un tirage aléatoire équivalent.

\begin{table}[H]
\centering
\caption{Description des variables d'un trade, illustrées sur un trade réel
de la base (achat du titre \texttt{A} par la stratégie \texttt{rsi\_classic}).}
\label{tab:dico-trades}
\renewcommand{\arraystretch}{1.3}
\small
\begin{tabular}{@{}>{\raggedright\arraybackslash}p{2.4cm}>{\raggedright\arraybackslash}p{2.6cm}>{\raggedright\arraybackslash}p{5.6cm}l@{}}
\toprule
Variable & Nature & Signification & Exemple \\
\midrule
\texttt{ticker}       & qualitative (501 mod.) & le titre traité (une même action revient dans plusieurs trades) & \texttt{A} \\
\texttt{strategy}     & qualitative (11 mod.)  & la stratégie ayant produit le trade & \texttt{rsi\_classic} \\
\texttt{sector}       & qualitative (11 mod.)  & le secteur d'activité du titre & Santé \\
\texttt{entry\_date}  & date                   & date d'achat & 2011-08-11 \\
\texttt{exit\_date}   & date                   & date de revente & 2012-01-20 \\
\texttt{holding\_days}& quantitative           & durée de détention, en jours & $162$ \\
\texttt{regime\_entry}& qualitative            & régime du marché à l'entrée & baissier \\
\texttt{vol\_entry}   & quantitative           & volatilité du titre à l'entrée & $0{,}043$ \\
\texttt{rsi\_entry}   & quantitative           & niveau de l'indicateur RSI à l'entrée & $30{,}1$ \\
\texttt{dist\_ma200}  & quantitative           & écart à la moyenne mobile sur 200 jours & $-0{,}210$ \\
\texttt{return\_net}  & quantitative           & rendement net réalisé par la stratégie & $+0{,}187$ \\
\texttt{rand\_return} & quantitative           & rendement d'une entrée au hasard équivalente & $+0{,}078$ \\
\texttt{edge}         & quantitative           & \textbf{avantage} = \texttt{return\_net} $-$ \texttt{rand\_return} & $+0{,}109$ \\
\texttt{win}          & binaire                & trade gagnant ($1$) ou perdant ($0$) & $1$ \\
\bottomrule
\end{tabular}
\end{table}

\section{Les rendements sectoriels}
\label{sec:donnees-secteurs}


Pour étudier la dépendance des données, on a besoin de données par secteur.
Un secteur regroupe les sociétés qui exercent le même type d'activité et qui
réagissent donc aux mêmes événements. Le S\&P~500 en compte onze, parmi
lesquels la technologie, la santé, l'énergie, ou encore les services aux
collectivités qui rassemblent les fournisseurs d'électricité, d'eau et de gaz.
Une hausse du prix du pétrole touchera ainsi toutes les sociétés du secteur de
l'énergie en même temps.

Sur chacun de ces secteurs on établit une courbe du rendement journalier,
obtenu en moyennant chaque jour les rendements des titres qui le composent.
Ces données seront utilisées d'une part dans le canal~2 du
paragraphe~\ref{sec:dep-canal2}, pour mettre en évidence le facteur
marché, c'est-à-dire la part des mouvements qui est commune à tous les
secteurs ; et d'autre part au paragraphe~\ref{sec:b1-efficience}, pour
examiner si le marché est prévisible à partir de son propre passé.



\section{La structure de dépendance, mesurée}
\label{sec:dep-mesuree}


Les tests que nous utiliserons au chapitre~\ref{chap:juger} ne fonctionnent que
si les données sont indépendantes.
Ici les données sont dépendantes par titre, par date et par durée.
Le réflexe du sondage impose de décrire son plan d'échantillonnage avant
d'estimer quoi que ce soit, c'est pourquoi on mesure ces trois dépendances
maintenant, avant les tests. On les examine l'une après l'autre, chacune
formant ce que nous appellerons un canal.

% --- Canal 1 : les grappes par titre ---------------------------------
\subsection{Canal 1 --- la dépendance de titre}
\label{sec:dep-canal1}


Deux trades d'un même titre sont liés par la gestion de la société
ainsi que par l'image de la société à l'extérieur.
Ces deux trades subissent les mêmes chocs propres à la société (résultats,
décisions de direction, litiges) indépendamment du moment où ils sont
ouverts.
On appellera cet effet la dépendance de titre : nos données constituent
alors un tirage par grappes, où chaque grappe est un titre.

Reste à mesurer l'ampleur de cette ressemblance. Pour cela on va définir
l'estimateur suivant, appelé corrélation intra-grappe, qui repose sur la
décomposition de la variance étudiée en analyse de la
variance~\cite{scheffe1959} :
\[
  \hat\rho \;=\; \frac{\mathrm{MSB}-\mathrm{MSW}}
                      {\mathrm{MSB} + (n_0 - 1)\,\mathrm{MSW}},
\]
où MSB (\emph{Mean Square Between}) est la variation entre les
titres, donc inter-grappes, c'est-à-dire la moyenne des écarts d'\edgevar{}
d'un titre à l'autre ; et MSW (\emph{Mean Square Within}) est la variation
à l'intérieur d'un titre, donc intra-grappe, c'est-à-dire la moyenne
des écarts d'\edgevar{} entre les trades d'un même titre.
Deux cas se présentent alors. Si les trades d'un même titre ne se ressemblent
pas plus que ceux de titres différents, la variation est identique à
l'intérieur des titres et entre eux, et l'on obtient MSB~$\approx$~MSW. S'ils
se ressemblent (dépendance intra-titre), la variation à l'intérieur d'un titre
s'effondre tandis que la variation entre les titres subsiste, et l'on obtient
MSB~$>$~MSW.

On pourrait prendre un simple quotient de MSB et MSW et lorsqu'il vaut~1,
cela signifierait l'indépendance. Mais un écart à~1, qu'il soit de
$0{,}1$, de~3 ou de~100, ne se rattache à aucune échelle. On ne saurait
dire s'il traduit une dépendance forte ou faible.
Pour obtenir une réponse du type «~13{,}4~\% du résultat est dû à une
dépendance de grappe~», on préfère la formule vue précédemment, qui ramène
la mesure à une valeur comprise entre~0 et~1.

Le terme $n_0$ correspond à la taille moyenne des grappes, c'est-à-dire au
nombre de trades par titre, corrigée de l'inégalité des tailles de grappes.
Il pondère la variation intra-titre au dénominateur, ce qui ramène la mesure
entre~0 et~1 : plus les grappes contiennent de trades, plus $n_0$ est grand,
et plus il faut de contraste entre les titres pour conclure à une dépendance.

Appliqué à nos données, cet estimateur, reporté au
tableau~\ref{tab:dep-mesures} en fin de section, donne un $\hat\rho$ quasi nul
pour la plupart des stratégies, mais atteignant $0{,}134$ pour
\texttt{rsi\_strict}.

La dépendance par le titre est donc négligeable pour toutes les stratégies,
excepté \texttt{rsi\_strict}. Un $\hat\rho$ faible ne garantit toutefois pas
l'indépendance des trades : il n'écarte que ce premier canal, les deux suivants
restant à examiner.


% --- Canal 2 : les vagues simultanées --------------------------------
\subsection{Canal 2 --- la dépendance de marché}
\label{sec:dep-canal2}


La deuxième dépendance que l'on va mesurer dans ce canal est la dépendance
du moment. Des trades ouverts au même moment peuvent être liés parce qu'ils
subissent la même journée de bourse : à une même période, différentes actions
changent de valeur conjointement, et les trades effectués sont alors liés par
le marché.

Cette dépendance ne s'observe pas directement sur les trades, mais sur le
marché lui-même. Si les secteurs bougent ensemble, alors des trades ouverts
au même moment, quel que soit leur secteur, subissent la même poussée. Pour
la mesurer, on utilise donc les données introduites au
point~\ref{sec:donnees-secteurs}, c'est-à-dire les rendements quotidiens par
secteur.
On calcule alors les corrélations des secteurs deux à deux, puis on
en fait la moyenne ; on obtient une corrélation de $0{,}70$.
Ce chiffre élevé signifie que les secteurs sont largement corrélés : lorsqu'un
secteur bouge, le marché a tendance à bouger dans le même sens, et
inversement, lorsque le marché suit une tendance, la plupart des secteurs la
suivent aussi.

On utilise ensuite la méthode de l'ACP~\cite{lebart2006} sur les onze secteurs, de
façon à mettre en évidence les plus grandes corrélations
(figure~\ref{fig:acp-marche}).
Cette méthode calcule les axes qui captent la plus grande
variance et les met en lumière.
En l'appliquant ici, on observe un premier axe qui capte
73~\% de la variance totale, le long duquel tous les secteurs
progressent ensemble : c'est le mouvement global du marché.
On a ensuite un deuxième axe, nettement moins important,
qui sépare les secteurs cycliques, réagissant fortement au
marché, de ceux qui sont défensifs, c'est-à-dire qui ne bougent
pas beaucoup quelle que soit la situation du marché. On pourrait
donner en exemple de secteur cyclique la technologie, qui dépend du pouvoir
d'achat ; et comme secteur défensif les services aux collectivités
(électricité, eau, gaz), dont la consommation reste constante, que le marché
soit haussier ou en crise.


\begin{figure}[H]
\centering
\begin{subfigure}[b]{0.46\linewidth}
\centering
\begin{tikzpicture}
\begin{axis}[
  width=\linewidth, height=5.4cm,
  ybar, bar width=6pt,
  xlabel={Composante}, ylabel={Variance expliquée (\%)},
  xlabel style={font=\small}, ylabel style={font=\small},
  tick label style={font=\footnotesize},
  xmin=0.3, xmax=11.7, ymin=0, ymax=80,
  xtick={1,2,3,4,5,6,7,8,9,10,11},
  ytick={0,20,40,60,80},
  ymajorgrids=true, grid style={dotted,gray!40},
  axis lines=left,
]
\addplot[fill=bleuunistra!75, draw=bleuunistra] coordinates {
  (1,73.0) (2,7.9) (3,5.1) (4,3.0) (5,2.3) (6,2.1)
  (7,1.9) (8,1.5) (9,1.3) (10,1.2) (11,0.6)
};
\end{axis}
\end{tikzpicture}
\caption{Variance expliquée par composante.}
\label{fig:acp-eboulis}
\end{subfigure}
\hfill
\begin{subfigure}[b]{0.50\linewidth}
\centering
\begin{tikzpicture}
\begin{axis}[
  width=\linewidth, height=5.4cm,
  xlabel={Axe 1 (73~\%)}, ylabel={Axe 2 (7,9~\%)},
  xlabel style={font=\small}, ylabel style={font=\small},
  tick label style={font=\footnotesize},
  xmin=0, xmax=0.40, ymin=-0.42, ymax=0.80,
  xtick={0,0.1,0.2,0.3,0.4},
  ytick={-0.4,-0.2,0,0.2,0.4,0.6,0.8},
  grid=both, grid style={dotted,gray!40},
  axis lines=left,
  scatter/classes={
    d={mark=*,bleuunistra,mark size=1.7pt},
    c={mark=triangle*,black!65,mark size=2.2pt}},
]
% Secteurs défensifs (disques) et cycliques (triangles), tracés séparément
% pour pouvoir placer chaque étiquette sans chevauchement.
\addplot[only marks, mark=*, bleuunistra, mark size=1.7pt] coordinates {
  (0.277,0.433) (0.307,0.023) (0.299,0.316) (0.244,0.674)
};
\addplot[only marks, mark=triangle*, black!65, mark size=2.2pt] coordinates {
  (0.310,-0.183) (0.321,-0.185) (0.255,-0.227) (0.328,-0.109)
  (0.335,-0.129) (0.304,-0.303) (0.321,-0.128)
};
\draw[dashed,gray!70] (axis cs:0,0) -- (axis cs:0.40,0);
% étiquettes placées à la main
\node[font=\tiny,anchor=east,black!75] at (axis cs:0.244,0.674) {Serv.\,coll.~};
\node[font=\tiny,anchor=east,black!75] at (axis cs:0.277,0.433) {Cons.\,base~};
\node[font=\tiny,anchor=west,black!75] at (axis cs:0.299,0.316) {~Immobilier};
\node[font=\tiny,anchor=west,black!75] at (axis cs:0.307,0.023) {~Santé};
\node[font=\tiny,anchor=east,black!75] at (axis cs:0.328,-0.109) {Finance~};
\node[font=\tiny,anchor=west,black!75] at (axis cs:0.335,-0.129) {~Industrie};
\node[font=\tiny,anchor=east,black!75] at (axis cs:0.321,-0.128) {Matér.~};
\node[font=\tiny,anchor=east,black!75] at (axis cs:0.310,-0.183) {Com.~};
\node[font=\tiny,anchor=west,black!75] at (axis cs:0.321,-0.185) {~Cons.\,disc.};
\node[font=\tiny,anchor=east,black!75] at (axis cs:0.255,-0.227) {Énergie~};
\node[font=\tiny,anchor=west,black!75] at (axis cs:0.304,-0.303) {~Techno.};
\end{axis}
\end{tikzpicture}
\caption{Position des secteurs sur les deux premiers axes.}
\label{fig:acp-nuage}
\end{subfigure}
\caption{ACP des onze rendements sectoriels. Sur le nuage de droite, les
triangles désignent les secteurs cycliques et les disques les secteurs
défensifs.}
\label{fig:acp-marche}
\end{figure}

Sur l'histogramme de gauche, la première barre écrase toutes les autres,
elle capte à elle seule 73~\% de la variance, contre moins de 8~\% pour la
deuxième. Sur le nuage de droite, tous les secteurs se situent du même côté
de l'axe horizontal (la première composante), ce qui confirme qu'ils
progressent ensemble ; le second axe (vertical) les répartit entre secteurs
défensifs et cycliques.

Cette cause commune n'est donc pas une hypothèse mais une mesure : 73~\% des
mouvements sont portés par une force unique. Deux trades ouverts au même
moment subissent largement le même sort, ce qui les rend dépendants du marché.

% --- Canal 3 : les chevauchements dans le temps ----------------------
\subsection{Canal 3 --- la dépendance temporelle}
\label{sec:dep-canal3}


Dans la suite du rapport, on va avoir besoin d'agréger les trades par mois.
Or les trades durent : en moyenne 123~jours, avec une médiane de 53~jours,
et plus de 78~\% d'entre eux dépassent un mois. Un même trade, encore ouvert,
est donc compté dans plusieurs mois consécutifs ; deux mois voisins partagent
ainsi des trades communs et se ressemblent d'un mois à l'autre.
Même agrégées par mois, les observations ne sont donc pas indépendantes.
C'est pourquoi on va mesurer la corrélation d'un mois à l'autre.

Pour cela, on utilise l'autocorrélation d'ordre~1, notée
$\hat\rho(1)$ : on décale la série mensuelle d'un mois et on corrèle la série
avec elle-même ainsi décalée. Elle mesure donc à quel point un mois ressemble
au précédent. On obtient les valeurs suivantes.
\begin{table}[H]
\centering
\caption{Les dépendances mesurées, stratégie par stratégie :
corrélation intra-titre $\hat\rho$ (canal~1) et autocorrélation
mensuelle $\hat\rho(1)$ (canal~3).}
\label{tab:dep-mesures}
\small
\begin{tabular}{@{}lrr@{}}
\toprule
Stratégie & $\hat\rho$ intra-titre & $\hat\rho(1)$ mensuelle \\
\midrule
\texttt{oracle} (témoin) & $0{,}053$ & $0{,}18$ \\
\midrule
\texttt{db\_bottom}    & $0{,}000$ & $0{,}36$ \\
\texttt{dt\_top}       & $0{,}000$ & $0{,}38$ \\
\texttt{hs\_classic}   & $0{,}087$ & $0{,}15$ \\
\texttt{hs\_inverse}   & $0{,}001$ & $0{,}31$ \\
\texttt{ma\_crossover} & $0{,}003$ & $0{,}56$ \\
\texttt{rsi\_classic}  & $0{,}000$ & $0{,}23$ \\
\texttt{rsi\_strict}   & $0{,}134$ & $0{,}03$ \\
\texttt{rsi\_trend}    & $0{,}010$ & $0{,}20$ \\
\texttt{sr\_breakdown} & $0{,}051$ & $0{,}11$ \\
\texttt{sr\_breakout}  & $0{,}004$ & $0{,}05$ \\
\bottomrule
\end{tabular}
\end{table}

Le $\hat\rho(1)$ est de $0{,}56$ pour \texttt{ma\_crossover}, sa valeur la plus
forte, tandis qu'il est faible pour \texttt{rsi\_strict} ($0{,}03$).
La dépendance temporelle est a priori forte pour les stratégies
à trades longs, faible pour celles à trades courts. Attention :
ces valeurs élevées ne traduisent pas une prévisibilité du
marché, mais le chevauchement mécanique des trades qui durent.

% --- Bilan des trois canaux ---
Les trois canaux de dépendance sont désormais mesurés. Aucun n'est
systématiquement élevé. La dépendance de titre reste souvent proche de zéro,
et même l'autocorrélation temporelle s'efface pour les stratégies à trades
courts. Mais chacun se manifeste fortement pour certaines stratégies, la
dépendance de titre pour \texttt{rsi\_strict}, la dépendance de marché pour
toutes, la dépendance temporelle pour \texttt{ma\_crossover} ; et le
chapitre~\ref{chap:juger} montrera qu'ils suffisent à renverser des verdicts.

On peut maintenant juger les stratégies en sachant exactement ce qui menace
les tests.
